
This work measures test accuracy variance arising from random seed initialization in simple multilayer perceptrons trained under deterministic conditions. No prior work has established a validated classical variance baseline for neural networks trained on MNIST with $N\geq 30$ independent random seeds and stability analysis. We trained 1-layer and 2-layer MLPs on MNIST and Fashion-MNIST datasets, using N=30 seeds per condition under full determinism (fixed random seeds, deterministic algorithms, fixed batch order). Variance measurements show task-dependent scaling: Fashion-MNIST exhibits variance of 0.35\% (1-layer) and 0.59\% (2-layer), while MNIST shows 0.04\% (1-layer) and 0.06\% (2-layer). Statistical tests confirm non-zero variance in all conditions (p < 0.05). The causal mechanism was validated through hypothesis decomposition: different seeds produce independent initial weights (mean pairwise distance 9.6-16.2, p < 0.000001), which lead to different final model states after training (final distances 22.7-27.3, coefficient of variation 2-3\% for final loss). Bootstrap stability analysis revealed that N=30 provides sufficient statistical power for variance detection (p < 0.05) but not for precise estimation (confidence interval widths 93-110\% vs. 50\% threshold). This finding refines existing sample size guidelines for neural network variance measurement. The work provides measurement infrastructure for uncertainty quantification research and identifies task-dependency constraints for variance-based methods.
